\documentclass[10pt,a4paper]{article}
\usepackage[T1]{fontenc}
\usepackage[margin=2.1cm]{geometry}
\usepackage{times}
\usepackage{booktabs}
\usepackage{amsmath}
\usepackage{enumitem}
\usepackage{xcolor}
\usepackage{pgfplots}
\pgfplotsset{compat=1.18}
\usepgfplotslibrary{groupplots}
\usetikzlibrary{patterns,calc}
\usepackage[hidelinks]{hyperref}
\setlist{itemsep=2pt,topsep=3pt}
\emergencystretch=2em
\title{\textbf{Context-Aware Alert Prioritization: test-set results of the locked run}\\[2pt]\large Separate results report, to be merged into Deliverable 3}
\date{4 November 2026}
\author{}
\begin{document}
\maketitle
\vspace{-2.5em}

\section*{Status}
Every number in this report comes from the one locked run of the pre-registered protocol (\path{provenance/ait/PREREGISTRATION.md}) on the four AIT-ADS \emph{test} scenarios fox, santos, wardbeck and wilson: tag \texttt{freeze-2026-11-04}, commit \texttt{b7857b0}, run on 3 November 2026 from 20:15 to 20:42 (+07), ten forest seeds, every check of the run script passed, no deviation (\path{provenance/ait/DEVIATIONS.md}). The outputs are in \path{provenance/ait/test/}; the figures and tables below are generated from them (Section~\ref{sec:repro}). What may be claimed was fixed before the run (pre-registration, Sections 4--6 and 9) and is recorded in \path{provenance/ait/CLAIMS.md}; this report uses only that wording. One analysis is post hoc and is marked as such wherever it appears (Section~\ref{sec:posthoc}).

\paragraph{Unit and measure.} The ten non-scan attack stage occurrences of the four scenarios (online password cracking, user change, sudo escalation), each judged by the rank of the first queue item that holds it; seeded systems enter by their mean rank over the ten seeds. \textbf{Systems:} B0 (Wazuh rule level), B0$_\mathrm{rank}$ (level rarity), B0$_\mathrm{host}$ (host rarity), B$_\mathrm{RBA}$ (risk-based alerting after Uetz et al., reimplemented and checked against their published values), B2 (the method), its ablations and B3 (B2 with entity risk accumulation).

\section{Confirmatory outcome}
Paired first-item ranks, two-sided Wilcoxon, Holm within each family (Table~\ref{tab:res-paired}).
\begin{itemize}
\item \textbf{B2 vs host rarity} (primary): B2 earlier on 6 of 10 stages, later on 4; $p = 0.846$. Not significant. AIT-ADS cannot separate temporal context from host rarity here; host rarity is strong on AIT-ADS because every attack lands on the scenario's least active host, not because of a general property of attacks.
\item \textbf{B2 vs level rarity} (primary): B2 earlier on 7 of 10, later on 3; $p = 0.049$, Holm $p = 0.098$. Not significant after correction (the smaller $p$ of the family had to be below 0.025).
\item \textbf{B2 vs risk-based alerting} (published-work family): B$_\mathrm{RBA}$ earlier on 9 of 10 stages; $p = 0.006$. \textbf{Risk-based alerting ranks attack stages earlier than the method.}
\item \textbf{Without reconnaissance} (Section 5a): B$_\mathrm{RBA}$ earlier on 7 of 10; $p = 0.131$. No difference between the method and risk-based alerting was found.
\end{itemize}
Calibration split: only wilson's two stages lie on a host calibrated mainly against its own baseline ($n_h \ge n_0$); on them B2 and host rarity split 1 to 1, too few to say anything. The other eight lie on hosts dominated by the pooled distribution.

\section{Figures}

\begin{figure}[h]
\centering
% generated by tools/make_test_figures.py from provenance/ait/test/{test_recorded,test_norecon}.json; do not edit by hand
\providecolor{csB2}{HTML}{2A78D6}
\providecolor{csRBA}{HTML}{EB6834}
\providecolor{csHost}{HTML}{1BAF7A}
\providecolor{csRank}{HTML}{EDA100}
\providecolor{csInk}{HTML}{52514E}
\providecolor{csGrid}{HTML}{DDDCD7}
\begin{tikzpicture}
\begin{axis}[width=\linewidth, height=5.4cm, ybar=1pt, bar width=9pt, axis lines*=left, tick align=outside, tick style={csInk}, axis line style={csInk}, ymajorgrids, grid style={csGrid, line width=0.3pt}, label style={font=\small}, tick label style={font=\small},
  ymin=0, ymax=10, ytick={0,2,4,6,8,10}, ylabel={stages first met in the top 20 (of 10)},
  xtick={0,...,6}, xticklabels={B0,B0$_\mathrm{rank}$,B0$_\mathrm{host}$,\textbf{B2},B2$_\mathrm{pooled}$,B3,B$_\mathrm{RBA}$}, x tick label style={font=\small}, enlarge x limits=0.08,
  area legend, legend style={draw=none, font=\small, at={(0.02,0.98)}, anchor=north west, fill=none}, legend cell align=left]
\addplot[fill=csInk, draw=csInk] coordinates {(0,0.000) (1,2.000) (2,7.000) (3,6.400) (4,8.000) (5,6.000) (6,10.000)};
\addplot[fill=white, draw=csInk, pattern=north east lines, pattern color=csInk] coordinates {(0,0.000) (1,3.000) (2,7.000) (3,7.500) (4,8.000) (5,7.100) (6,9.000)};
\legend{recorded (headline), without reconnaissance}
\end{axis}
\end{tikzpicture}

\caption{Attack stages found early. Non-scan stages (of 10) whose first queue item is in the top 20, per system, summed over the four test scenarios (seeded systems: mean over ten seeds). Solid: recorded baselines (headline); hatched: scan alerts removed from the stream. B0, Wazuh's own rule-level ordering, places none of the ten in the top 20 in either run (empty bars). B2 against B0 is a secondary, exploratory comparison (B2 earlier on 9 of 10 stages in both runs); B2$_\mathrm{pooled}$ is a pre-registered ablation, not the method.}
\label{fig:coverage}
\end{figure}

Compared with Wazuh's rule-level ordering, the layer placed 9 of 10 attack stages earlier, in both runs (exploratory; Holm $p = 0.107$ recorded, $0.043$ without reconnaissance). Risk-based alerting places all ten stages in the top 20 of the recorded run and nine without reconnaissance.

\begin{figure}[h]
\centering
% generated by tools/make_test_figures.py from provenance/ait/test/{test_recorded,test_norecon}.json; do not edit by hand
\providecolor{csB2}{HTML}{2A78D6}
\providecolor{csRBA}{HTML}{EB6834}
\providecolor{csHost}{HTML}{1BAF7A}
\providecolor{csRank}{HTML}{EDA100}
\providecolor{csInk}{HTML}{52514E}
\providecolor{csGrid}{HTML}{DDDCD7}
\begin{tikzpicture}
\begin{axis}[width=\linewidth, height=5.4cm, ybar=1pt, bar width=11pt, axis lines*=left, tick align=outside, tick style={csInk}, axis line style={csInk}, ymajorgrids, grid style={csGrid, line width=0.3pt}, label style={font=\small}, tick label style={font=\small},
  ymin=0, ymax=1.05, ytick={0,0.25,0.5,0.75,1}, ylabel={within-host AUC},
  xtick={0,...,3}, xticklabels={fox,santos,wardbeck,wilson}, enlarge x limits=0.15,
  legend style={draw=none, font=\small, at={(0.5,1.02)}, anchor=south, fill=none, legend columns=2, /tikz/every even column/.append style={column sep=8pt}},
  area legend, legend cell align=left]
\addplot[fill=csRank, draw=white, line width=0.6pt] coordinates {(0,0.486) (1,0.517) (2,0.540) (3,0.544)};
\addplot[fill=csB2, draw=white, line width=0.6pt] coordinates {(0,0.583) (1,0.953) (2,0.887) (3,0.940)};
\addplot[fill=csRBA, draw=white, line width=0.6pt] coordinates {(0,0.902) (1,0.988) (2,0.961) (3,0.882)};
\draw[csHost, line width=1.2pt, dash pattern=on 4pt off 2pt] (axis cs:-0.5,0.5) -- (axis cs:3.5,0.5);
\addlegendimage{line legend, csHost, line width=1.2pt, dash pattern=on 4pt off 2pt}
\legend{level rarity B0$_\mathrm{rank}$, B2, B$_\mathrm{RBA}$, host rarity (0.5)}
\end{axis}
\end{tikzpicture}

\caption{Separation within a host. For each host carrying at least five non-scan attack and five benign alerts (one host per scenario), the AUC of its attack alerts against its own benign alerts, recorded run. Host rarity gives every alert of a host the same score and is 0.5 on every host by construction (dashed line). For B2 the AUC is computed on the window score $S_\mathrm{seq}$, mean over seeds. Descriptive secondary measure on four hosts.}
\label{fig:withinhost}
\end{figure}

Within a host, the window score separates attack from benign alerts on all four eligible hosts (AUC 0.58 to 0.95), where host rarity and level rarity are at chance (0.50 and 0.49--0.54); risk-based alerting does so better on three of the four (pre-registered descriptive comparison, Section 9.1: B2 higher on 1 of 4 hosts, verdict mixed).

\begin{figure}[h]
\centering
% generated by tools/make_test_figures.py from provenance/ait/test/{test_recorded,test_norecon}.json; do not edit by hand
\providecolor{csB2}{HTML}{2A78D6}
\providecolor{csRBA}{HTML}{EB6834}
\providecolor{csHost}{HTML}{1BAF7A}
\providecolor{csRank}{HTML}{EDA100}
\providecolor{csInk}{HTML}{52514E}
\providecolor{csGrid}{HTML}{DDDCD7}
\begin{tikzpicture}
\begin{axis}[width=\linewidth, height=5.4cm, ybar=1pt, bar width=9pt, axis lines*=left, tick align=outside, tick style={csInk}, axis line style={csInk}, ymajorgrids, grid style={csGrid, line width=0.3pt}, label style={font=\small}, tick label style={font=\small},
  ymin=0, ymax=0.8, ytick={0,0.2,0.4,0.6,0.8}, ylabel={$\alpha$-nDCG@20},
  xtick={0,...,6}, xticklabels={B0,B0$_\mathrm{rank}$,B0$_\mathrm{host}$,\textbf{B2},B2$_\mathrm{pooled}$,B3,B$_\mathrm{RBA}$}, x tick label style={font=\small}, enlarge x limits=0.08,
  area legend, legend style={draw=none, font=\small, at={(0.02,0.98)}, anchor=north west, fill=none}, legend cell align=left]
\addplot[fill=csInk, draw=csInk] coordinates {(0,0.233) (1,0.289) (2,0.332) (3,0.327) (4,0.429) (5,0.412) (6,0.676)};
\addplot[fill=white, draw=csInk, pattern=north east lines, pattern color=csInk] coordinates {(0,0.000) (1,0.113) (2,0.330) (3,0.210) (4,0.357) (5,0.197) (6,0.299)};
\legend{recorded (headline), without reconnaissance}
\end{axis}
\end{tikzpicture}

\caption{Dependence on reconnaissance. Graded $\alpha$-nDCG@20 per system (mean over scenarios; seeded systems over ten seeds), with recorded baselines (solid) and with the service, directory and WordPress scan alerts removed from the stream (hatched). The pre-registered test of B2 against B$_\mathrm{RBA}$ without reconnaissance is not significant (B$_\mathrm{RBA}$ earlier on 7 of 10 stages, $p = 0.131$); the changes in the aggregate measure are descriptive and do not establish robustness.}
\label{fig:recon}
\end{figure}

Removing reconnaissance lowered risk-based alerting's $\alpha$-nDCG@20 from 0.676 to 0.299 ($-56\%$) and its MRR from 0.425 to 0.106, and B2's $\alpha$-nDCG@20 from 0.327 to 0.210 ($-36\%$) and MRR from 0.134 to 0.114; level rarity fell by 61\% and host rarity by 1\%. The paired test above found no difference between B2 and risk-based alerting in this condition.

\begin{figure}[h]
\centering
% generated by tools/make_test_figures.py from provenance/ait/test/{test_recorded,test_norecon}.json; do not edit by hand
\providecolor{csB2}{HTML}{2A78D6}
\providecolor{csRBA}{HTML}{EB6834}
\providecolor{csHost}{HTML}{1BAF7A}
\providecolor{csRank}{HTML}{EDA100}
\providecolor{csInk}{HTML}{52514E}
\providecolor{csGrid}{HTML}{DDDCD7}
\begin{tikzpicture}
\begin{groupplot}[group style={group size=2 by 1, horizontal sep=0.35cm}, scale only axis, width=0.36\linewidth, height=6.4cm,
  xmode=log, xmin=0.8, xmax=110, xtick={1,2,5,10,20,50,100}, xticklabels={1,2,5,10,20,50,100},
  xlabel={rank of the first item (log; lower = earlier)}, ymin=0.4, ymax=10.6,
  axis lines*=left, tick align=outside, tick style={csInk}, axis line style={csInk},
  xmajorgrids, ymajorgrids, grid style={csGrid, line width=0.3pt}, label style={font=\small},
  tick label style={font=\small}, title style={font=\small},
  legend style={draw=none, font=\small, fill=none, legend columns=2, /tikz/every even column/.append style={column sep=8pt}}]
\nextgroupplot[title={recorded (headline)}, ytick={1,...,10}, yticklabels={wilson: sudo,wilson: user change,wardbeck: sudo,wardbeck: user change,santos: sudo,santos: user change,santos: cracking,fox: sudo,fox: user change,fox: cracking}, legend to name=firstranklegend,]
\addplot[only marks, color=csHost, mark=*, mark size=2.4pt] coordinates {(30.0,10) (30.0,9) (30.0,8) (10.0,7) (10.0,6) (10.0,5) (17.0,4) (17.0,3) (17.5,2) (17.5,1)}; \addlegendentry{host rarity B0$_\mathrm{host}$}
\addplot[only marks, color=csRBA, mark=square*, mark size=2.2pt] coordinates {(7.0,10) (2.0,9) (2.0,8) (8.0,7) (3.0,6) (1.0,5) (2.0,4) (2.0,3) (3.0,2) (3.0,1)}; \addlegendentry{risk-based alerting B$_\mathrm{RBA}$}
\addplot[only marks, color=csB2, mark=triangle*, mark size=2.8pt] coordinates {(10.6,10) (10.9,9) (65.4,8) (6.3,7) (4.0,6) (10.8,5) (13.6,4) (26.6,3) (7.2,2) (45.7,1)}; \addlegendentry{B2}
\addplot[only marks, color=csB2, mark=diamond, mark size=2.8pt, mark options={line width=1pt, fill=white}] coordinates {(6.0,10) (6.8,9) (49.1,8) (6.5,7) (4.2,6) (8.4,5) (2.2,4) (5.5,3) (3.8,2) (37.9,1)}; \addlegendentry{B2$_\mathrm{pooled}$ (ablation)}
\nextgroupplot[title={without reconnaissance}, ytick={1,...,10}, yticklabels={},]
\addplot[only marks, color=csHost, mark=*, mark size=2.4pt] coordinates {(24.0,10) (24.0,9) (24.0,8) (6.0,7) (6.0,6) (6.0,5) (12.5,4) (12.5,3) (12.0,2) (12.0,1)};
\addplot[only marks, color=csRBA, mark=square*, mark size=2.2pt] coordinates {(13.0,10) (9.0,9) (8.0,8) (33.0,7) (10.0,6) (8.0,5) (8.0,4) (8.0,3) (9.0,2) (9.0,1)};
\addplot[only marks, color=csB2, mark=triangle*, mark size=2.8pt] coordinates {(11.2,10) (18.1,9) (52.0,8) (2.3,7) (8.4,6) (9.9,5) (12.2,4) (19.5,3) (16.7,2) (27.9,1)};
\addplot[only marks, color=csB2, mark=diamond, mark size=2.8pt, mark options={line width=1pt, fill=white}] coordinates {(6.5,10) (8.8,9) (34.8,8) (2.8,7) (3.8,6) (4.7,5) (1.4,4) (3.6,3) (3.5,2) (23.4,1)};
\end{groupplot}
\node[anchor=north] at ($(group c1r1.south east)!0.5!(group c2r1.south west) + (0,-1.1cm)$) {\pgfplotslegendfromname{firstranklegend}};
\end{tikzpicture}

\caption{Where each stage is first met. Rank of the first queue item holding each non-scan stage occurrence (log scale; lower is earlier), recorded run (left) and without reconnaissance (right); seeded systems: mean over ten seeds. B2$_\mathrm{pooled}$ (hollow) is a pre-registered ablation; its comparison with B$_\mathrm{RBA}$ in this figure is post hoc (Section~\ref{sec:posthoc}).}
\label{fig:firstrank}
\end{figure}

B2 reaches password cracking and the user change early (ranks 4 to 14 in the recorded run) but sudo escalation late (ranks 65.4 in fox, 10.8 santos, 26.6 wardbeck, 45.7 wilson), where risk-based alerting reaches it at ranks 1 to 3.

\clearpage
\section{Secondary comparisons (exploratory)}
All twelve comparisons of the secondary family are listed with counts and $p$-values in Table~\ref{tab:res-paired}; none is stated as a confirmed improvement.
\begin{itemize}
\item \textbf{The forest does not earn its place:} B2 against the robust $z$-score B2$'$, 7 to 3 ($p = 0.695$) recorded and 3 to 7 without reconnaissance. The robust $z$-score is reported as an equal alternative, and the result is credited to the window representation rather than to the Isolation Forest.
\item \textbf{Per-host calibration does not improve ranking:} pooled calibration (B2$_\mathrm{pooled}$) placed 8 of 10 stages earlier than the specified per-host calibration in the recorded run (Holm $p = 0.107$) and 9 of 10 without reconnaissance (Holm $p = 0.043$). Its within-host AUC equals B2's: calibration changes only how hosts are ordered against each other.
\item The size correction $\pi(g)$ is the one component whose comparison has a Holm $p$ below 0.05 in the recorded run (B2 earlier than B2$_\mathrm{max}$ on 10 of 10, Holm $p = 0.023$); the timing features only without reconnaissance (10 of 10, Holm $p = 0.023$; 3 to 7 in the recorded run).
\item The window against the novelty flag alone (B2$_\mathrm{nov}$): 9 to 1 recorded ($p = 0.014$, Holm $0.123$), 5 to 5 without reconnaissance.
\item Entity accumulation (B3) does not change B2's first-item ranks (2 earlier, 5 later, 3 tied); against B$_\mathrm{RBA}$ it is earlier on 1 of 10 stages recorded and 3 of 10 without reconnaissance (descriptive, Section 9.3).
\end{itemize}

\section{Santos: recorded and cleaned baselines}
The headline uses santos as recorded, with its designed exfiltration in the baseline; Table~\ref{tab:res-cleaned} reports it beside the run with those alerts removed from the stream before the first attack step. On this one scenario B2 ranked better with the recorded baseline ($\alpha$-nDCG@20 0.499 recorded, 0.390 cleaned; first items of user change 4.0 and 9.2, sudo escalation 10.8 and 14.3, cracking 6.3 and 12.8). This is evidence from one scenario only and is not a general claim of tolerance to contamination.

\section{Post-hoc observations}
\label{sec:posthoc}
Not part of the protocol; listed in \texttt{DEVIATIONS.md} (P1); descriptive, without tests. They read only the frozen outputs (\path{tools/ait_posthoc.py}, \path{provenance/ait/posthoc_2026-11-04.txt}).
\begin{itemize}
\item \textbf{The pooled-calibration ablation against risk-based alerting.} B2$_\mathrm{pooled}$ reached the first item earlier than B$_\mathrm{RBA}$ on 2 of 10 stages in the recorded run and on 8 of 10 without reconnaissance ($\alpha$-nDCG@20 0.357 against 0.299, MRR 0.276 against 0.106), and earlier than host rarity on 8 of 10 in both runs. This suggests pooled calibration as the better design choice; the comparison was not pre-registered, B2$_\mathrm{pooled}$ is not the method, and the observation needs confirmation on new data. Part of its gain may be host rarity re-entering the score, since pooled calibration rates a window on a rarely active host as more unusual.
\item \textbf{Calibration of one host.} The per-host calibration check expects about 1\% of benign evaluation windows above $S_\mathrm{seq} = 99$; on \texttt{internal\_share} it found 67\% (fox), 100\% (santos) and 27\% (wilson). That host's behaviour after the split departs from its baseline, and its benign groups occupy queue positions above the late stages. This is an observation, not a demonstrated cause of the late sudo escalation.
\end{itemize}

\section{How to reproduce}
\label{sec:repro}
From the repository root, with the environment of \path{tools/requirements-ait.txt}. The locked run itself is
\texttt{bash tools/ait\_locked\_run.sh freeze-2026-11-04} at that tag (procedure in \path{docs/LOCKED_RUN_CHECKLIST.md});
it refuses to run while \path{provenance/ait/test/} exists, and the committed outputs are the record. Figures, tables
(identical to the run's own \path{results_tables.tex} apart from the corrected caption of Table~\ref{tab:res-cleaned})
and post-hoc counts:
{\small
\begin{verbatim}
python3 tools/make_test_figures.py tools provenance/ait/test References/results_test
python3 tools/make_result_tables.py tools provenance/ait/test/test_recorded.json \
    References/results_test/results_tables.tex \
    --norecon provenance/ait/test/test_norecon.json \
    --cleaned provenance/ait/test/test_santos_cleaned.json
python3 tools/ait_posthoc.py tools provenance/ait/test
\end{verbatim}}

\clearpage
% generated by tools/make_result_tables.py from test_recorded.json and test_norecon.json and test_santos_cleaned.json; do not edit by hand
\begin{table*}[t]
\caption{All systems on fox, santos, wardbeck, wilson: graded $\alpha$-nDCG@$k$ (mean over scenarios; seeded systems over 10 seeds, $\pm$ seed SD of @20), MRR of non-scan stages, non-scan stages first met in the top 20, macro-AUC and within-host AUC}
\label{tab:res-main}
\centering\footnotesize
\begin{tabular}{@{}lccccccc@{}}
\toprule
\textbf{System} & \textbf{@10} & \textbf{@20} & \textbf{@50} & \textbf{MRR} & \textbf{Cov@20} & \textbf{macro-AUC} & \textbf{within-host AUC} \\
\midrule
B0 & 0.218 & 0.233 & 0.352 & 0.021 & 0.0/10 & 0.521 & 0.521 \\
B1 & 0.203 & 0.258 & 0.308 & 0.028 & 2.0/10 & 0.428 & 0.521 \\
B0$_\mathrm{rank}$ & 0.270 & 0.289 & 0.365 & 0.043 & 2.0/10 & 0.570 & 0.522 \\
B0$_\mathrm{host}$ & 0.206 & 0.332 & 0.431 & 0.062 & 7.0/10 & 0.943 & 0.500 \\
B$_\mathrm{RBA}$ & 0.611 & 0.676 & 0.692 & 0.425 & 10.0/10 & 0.981 & 0.933 \\
B2 & 0.269 & 0.327$\pm$0.040 & 0.405 & 0.134 & 6.4/10 & 0.919 & 0.841 \\
B3 & 0.368 & 0.412$\pm$0.031 & 0.491 & 0.108 & 6.0/10 & -- & -- \\
\midrule
B2$'$ & 0.231 & 0.306 & 0.392 & 0.086 & 8.0/10 & 0.781 & 0.840 \\
B2$_\mathrm{nov}$ & 0.293 & 0.357 & 0.436 & 0.044 & 5.0/10 & 0.707 & 0.695 \\
B2$_\mathrm{seq}$ & 0.253 & 0.310$\pm$0.048 & 0.403 & 0.130 & 5.9/10 & 0.919 & 0.841 \\
B2$_\mathrm{alert}$ & 0.218 & 0.259$\pm$0.028 & 0.358 & 0.118 & 4.7/10 & 0.900 & 0.701 \\
B2$_\mathrm{max}$ & 0.319 & 0.382$\pm$0.038 & 0.438 & 0.069 & 5.5/10 & -- & -- \\
B2$_\mathrm{pooled}$ & 0.357 & 0.429$\pm$0.017 & 0.490 & 0.207 & 8.0/10 & 0.974 & 0.841 \\
B2$_\mathrm{notime}$ & 0.290 & 0.341$\pm$0.046 & 0.425 & 0.157 & 5.7/10 & 0.928 & 0.852 \\
B2$_\mathrm{def}$ & 0.200 & 0.274$\pm$0.032 & 0.355 & 0.066 & 5.4/10 & -- & -- \\
B3$_\mathrm{bestgrp}$ & 0.271 & 0.331$\pm$0.042 & 0.410 & 0.135 & 6.4/10 & -- & -- \\
B3$_\mathrm{noabs}$ & 0.366 & 0.410$\pm$0.032 & 0.480 & 0.107 & 5.9/10 & -- & -- \\
\bottomrule
\end{tabular}
\end{table*}
\begin{table}[t]
\caption{Pre-registered paired comparisons of first-item ranks of non-scan stage occurrences: occurrences where the first system is earlier, later or tied; two-sided Wilcoxon $p$; Holm $p$ within each family}
\label{tab:res-paired}
\centering\footnotesize
\setlength{\tabcolsep}{3pt}
\begin{tabular}{@{}llccccc@{}}
\toprule
\textbf{Family} & \textbf{Comparison} & \textbf{earlier} & \textbf{later} & \textbf{tied} & $p$ & \textbf{Holm} $p$ \\
\midrule
primary & B2 vs B0$_\mathrm{host}$ & 6 & 4 & 0 & 0.846 & 0.846 \\
primary & B2 vs B0$_\mathrm{rank}$ & 7 & 3 & 0 & 0.049 & 0.098 \\
\midrule
published & B2 vs B$_\mathrm{RBA}$ & 1 & 9 & 0 & 0.006 & 0.006 \\
\midrule
secondary & B2 vs B0 & 9 & 1 & 0 & 0.010 & 0.107 \\
secondary & B2 vs B1 & 8 & 2 & 0 & 0.027 & 0.219 \\
secondary & B2 vs B2$'$ & 7 & 3 & 0 & 0.695 & 1.000 \\
secondary & B2 vs B2$_\mathrm{nov}$ & 9 & 1 & 0 & 0.014 & 0.123 \\
secondary & B2 vs B2$_\mathrm{seq}$ & 8 & 2 & 0 & 0.432 & 1.000 \\
secondary & B2 vs B2$_\mathrm{alert}$ & 6 & 4 & 0 & 0.160 & 0.801 \\
secondary & B2 vs B2$_\mathrm{max}$ & 10 & 0 & 0 & 0.002 & 0.023 \\
secondary & B2 vs B2$_\mathrm{pooled}$ & 2 & 8 & 0 & 0.010 & 0.107 \\
secondary & B2 vs B2$_\mathrm{notime}$ & 3 & 7 & 0 & 0.695 & 1.000 \\
secondary & B3 vs B2 & 2 & 5 & 3 & 0.656 & 1.000 \\
secondary & B3 vs B3$_\mathrm{bestgrp}$ & 0 & 6 & 4 & 0.031 & 0.219 \\
secondary & B3 vs B3$_\mathrm{noabs}$ & 6 & 0 & 4 & 0.031 & 0.219 \\
\bottomrule
\end{tabular}
\end{table}
\begin{table}[t]
\caption{First-item rank of every non-scan stage occurrence (seeded systems: mean over seeds; lower is earlier)}
\label{tab:res-stages}
\centering\scriptsize
\setlength{\tabcolsep}{3pt}
\begin{tabular}{@{}llrrrrr@{}}
\toprule
\textbf{Scenario} & \textbf{Stage} & \textbf{B0$_\mathrm{host}$} & \textbf{B0$_\mathrm{rank}$} & \textbf{B$_\mathrm{RBA}$} & \textbf{B2} & \textbf{B3} \\
\midrule
fox & user change & 30.0 & 83.0 & 2.0 & 10.9 & 11.0 \\
fox & sudo esc. & 30.0 & 83.0 & 2.0 & 65.4 & 61.4 \\
fox & password cracking & 30.0 & 59.5 & 7.0 & 10.6 & 10.7 \\
santos & user change & 10.0 & 44.5 & 3.0 & 4.0 & 4.0 \\
santos & sudo esc. & 10.0 & 7.5 & 1.0 & 10.8 & 10.8 \\
santos & password cracking & 10.0 & 44.5 & 8.0 & 6.3 & 6.3 \\
wardbeck & user change & 17.0 & 54.5 & 2.0 & 13.6 & 15.4 \\
wardbeck & sudo esc. & 17.0 & 7.5 & 2.0 & 26.6 & 27.8 \\
wilson & user change & 17.5 & 85.5 & 3.0 & 7.2 & 10.2 \\
wilson & sudo esc. & 17.5 & 28.0 & 3.0 & 45.7 & 43.9 \\
\bottomrule
\end{tabular}
\end{table}
\begin{table}[t]
\caption{No-reconnaissance condition: scan alerts removed from the stream; $\alpha$-nDCG@20 and MRR beside the main run, and the pre-registered paired test of B2 against B$_\mathrm{RBA}$}
\label{tab:res-norecon}
\centering\footnotesize
\begin{tabular}{@{}lcccc@{}}
\toprule
& \multicolumn{2}{c}{\textbf{main run}} & \multicolumn{2}{c}{\textbf{no recon.}} \\
\textbf{System} & @20 & MRR & @20 & MRR \\
\midrule
B0$_\mathrm{host}$ & 0.332 & 0.062 & 0.330 & 0.093 \\
B0$_\mathrm{rank}$ & 0.289 & 0.043 & 0.113 & 0.076 \\
B$_\mathrm{RBA}$ & 0.676 & 0.425 & 0.299 & 0.106 \\
B2 & 0.327 & 0.134 & 0.210 & 0.114 \\
B3 & 0.412 & 0.108 & 0.197 & 0.107 \\
\midrule
\multicolumn{5}{@{}p{0.9\columnwidth}@{}}{B2 vs B$_\mathrm{RBA}$ without reconnaissance: B2 earlier 3, later 7, tied 0; $p$ = 0.131} \\
\bottomrule
\end{tabular}
\end{table}
\begin{table}[t]
\caption{Contaminated scenario: recorded baseline (headline) and the stream with its labelled exfiltration alerts removed before the first attack step, that is, from training and from what the first day's evaluation windows look back on}
\label{tab:res-cleaned}
\centering\footnotesize
\begin{tabular}{@{}llcccc@{}}
\toprule
\textbf{Scenario} & \textbf{System} & \multicolumn{2}{c}{\textbf{recorded}} & \multicolumn{2}{c}{\textbf{cleaned}} \\
& & @20 & MRR & @20 & MRR \\
\midrule
santos & B$_\mathrm{RBA}$ & 0.725 & 0.486 & 0.718 & 0.492 \\
santos & B2 & 0.499 & 0.191 & 0.390 & 0.115 \\
santos & B3 & 0.499 & 0.191 & 0.465 & 0.092 \\
\bottomrule
\end{tabular}
\end{table}
\begin{table}[t]
\caption{Entity risk accumulation (B3), seed 0: incidents, those containing an attack alert, benign incidents, hours from the first attack step to the first attack notable, and the most stages in one incident}
\label{tab:res-b3}
\centering\footnotesize
\setlength{\tabcolsep}{4pt}
\begin{tabular}{@{}lccccc@{}}
\toprule
\textbf{Scenario} & \textbf{incid.} & \textbf{attack} & \textbf{benign} & \textbf{time (h)} & \textbf{stages} \\
\midrule
fox & 1 & 1 & 0 & 27.2 & 3 \\
santos & 0 & 0 & 0 & -- & 0 \\
wardbeck & 2 & 1 & 1 & 62.0 & 3 \\
wilson & 3 & 1 & 2 & 24.2 & 3 \\
\bottomrule
\end{tabular}
\end{table}


\end{document}
